\section{A second-order condition over the moduli field}\label{sec:second-order}

Normalize three distinct parameters to $0,1,\infty$. Choose a rational
$\alpha$ different from every finite normalized parameter, including
$0$ and $1$, and apply
$r\mapsto1/(r-\alpha)$. Every parameter is now finite and the first
three are distinct rational numbers. This operation and its inverse
are rational over $\Q$, so the field of all finite parameters remains
the moduli field $F=M(A)$. The unit circle coordinates have this same
field, since
\[
 a(t)=\left(\frac{1-t^2}{1+t^2},\frac{2t}{1+t^2}\right),\qquad
 t=\frac{a_2}{1+a_1}.
\]
The denominator in the inverse is $2/(1+t^2)>0$.
Section~\ref{sec:transport} justifies using this normalization for
certificate existence without a Ramsey-invariance input.

Write the distinct finite parameters as $t_1,\ldots,t_n$ and let
$k=\trdeg_\Q F$. Choose a transcendence basis $u_1,\ldots,u_k$.
The field $F$ is finite separable over $\Q(u_1,\ldots,u_k)$.
Each partial derivation extends uniquely to $F$: for an algebraic
element $\beta$ with minimal polynomial $f$, differentiating
$f(\beta)=0$ prescribes
\[
 D(\beta)=-\frac{(Df)(\beta)}{f'(\beta)}.
\]
Here $Df$ differentiates the coefficients. Separability makes the
denominator nonzero, and the prescription preserves the defining
ideal of the simple extension. Applying this to finite simple
extensions proves existence and uniqueness. It follows that
$\Der_\Q(F,F)$ has basis $D_1,\ldots,D_k$, where
$D_j(u_l)=\delta_{jl}$. Its common constants are precisely the
elements of $F$ algebraic over $\Q$: a transcendental element can
itself be included in a transcendence basis and assigned derivative
one.

Let
\begin{equation}\label{eq:velocity-space}
 V=\bigl\{(D(t_i))_{i=1}^n:D\in\Der_\Q(F,F)\bigr\},\qquad
 W=\bigl\{(q(t_i))_{i=1}^n:\deg q\le2\bigr\}.
\end{equation}
Since the $t_i$ generate $F$, the velocity map is injective, and
$\dim_F V=k$. Its first three entries vanish. A quadratic vanishing
at those three distinct parameters is zero, so
\begin{equation}\label{eq:transversality}
 V\cap W=0,\qquad \dim W=3.
\end{equation}

Set
\begin{equation}\label{eq:barycentric}
 H(T)=\prod_{i=1}^n(T-t_i),\qquad \gamma_i=H'(t_i)^{-1}.
\end{equation}
Lagrange interpolation, followed by comparison of the coefficient of
$T^{n-1}$, gives
\begin{equation}\label{eq:moments}
 \sum_i\gamma_i t_i^r=0\quad(0\le r\le n-2),\qquad
 \sum_i\gamma_i t_i^{n-1}=1.
\end{equation}

\begin{proposition}[The graph condition]\label{prop:graph}
Suppose $n\ge7$ and the circle set has a tensor certificate. There is
an $F$-linear map $Q:V\to W$ such that
\begin{equation}\label{eq:graph-relation}
 B_l(v,Qw)+B_l(w,Qv)=4B_l(v,w)
 \quad(v,w\in V,\ 0\le l\le n-6),
\end{equation}
where $B_l(a,b)=\sum_i\gamma_i t_i^l a_ib_i$.
Consequently
\begin{equation}\label{eq:graph-U}
 U=\{v-\tfrac14Qv:v\in V\}
\end{equation}
has dimension $k$, is totally isotropic for every $B_l$, and satisfies
$V\subset U+W$.
\end{proposition}
\begin{proof}
Normalize the certificate by \eqref{eq:full-normalization}. In the
truncated ring
\[
 \cR=F[\delta_1,\ldots,\delta_k]/(\delta_1,\ldots,\delta_k)^3
\]
put $D_\delta=\sum_j\delta_jD_j$, acting trivially on the $\delta_j$.
The map
\[
 \varphi(x)=x+D_\delta x+\tfrac12D_\delta^2x
\]
is a $\Q$-algebra homomorphism. This follows by applying Leibniz
twice and discarding terms of degree at least three. In particular
it respects inverses, since a nonzero constant term is a unit.
Commutation of the individual derivations is not needed for this
calculation. Map the left factor of $\cB_F$ identically to $F$ and
the right factor by $\varphi$.

Using $h(1,t)$, the image matrix represents a polynomial
\[
 \cG(t,u)=-2(t-u)^2+\cG_1(t,u)+\cG_2(t,u),
\]
of degree at most two in each variable, with $\cG_j$ homogeneous
of degree $j$ in $\delta$. The first-order node equations give
$\cG_1(t_i,t_i)=0$. Each coefficient on the diagonal is a polynomial
of degree at most four with at least seven roots; hence
$\cG_1(t,t)=0$. Polynomial division gives
\[
 \cG_1(t,u)=(u-t)N(t,u),\qquad
 N(t,t)=\sum_j a_j(t)\delta_j,\quad \deg a_j\le2.
\]
The quotient $N$ has degree at most one in each variable: multiplication
by $u-t$ raises the degree in either variable by one. Put
$q(t)=\cG_2(t,t)$, a polynomial in $t$ of degree at most four.
With $v_i=\sum_jD_j(t_i)\delta_j$, the second-order node equations
are
\begin{equation}\label{eq:quadratic-node}
 q(t_i)=2v_i^2-v_i\sum_j a_j(t_i)\delta_j.
\end{equation}
The quadratic displacement in $\varphi(t_i)$ contributes no term of
order two: the leading kernel has a double zero on the diagonal,
and $\cG_1$ already has order one in $\delta$.

For a velocity $v=\sum_j c_j(D_j(t_i))_i$, define
$Qv=(\sum_jc_j a_j(t_i))_i$. This is well-defined because the
velocity columns are independent, and it takes values in $W$.
Multiplying \eqref{eq:quadratic-node} by $\gamma_i t_i^l$ and
summing, \eqref{eq:moments} gives
$B_l(v,Qv)=2B_l(v,v)$ for $0\le l\le n-6$.
Polarization proves \eqref{eq:graph-relation}.
Also \eqref{eq:moments} makes $W$ totally isotropic for every
such $B_l$, because $l+4\le n-2$. Expanding
$B_l(v-Qv/4,w-Qw/4)$ now gives zero. The graph map is injective
by \eqref{eq:transversality}, so $\dim U=k$; its definition gives
$V\subset U+W$.
\end{proof}

There is a useful equivalent form of this condition. The
$(n-5)\times n$ matrix $(\gamma_i t_i^l)_{0\le l\le n-6}$ has
rank $n-5$ by the Vandermonde determinant. Its kernel is exactly
the five-dimensional space of values of polynomials of degree at
most four, by \eqref{eq:moments}. Thus
\begin{equation}\label{eq:product-property}
 u,w\in U\quad\Longrightarrow\quad
 (u_iw_i)_{i=1}^n\text{ is the value vector of a polynomial of degree }\le4.
\end{equation}
This formulation will allow restriction to eight nodes.
